% !TEX TS-program = lualatex

%%%
% src::
%     url = https://github.com/pfradin/luadraw/discussions/104#discussioncomment-14977791
%%%

\documentclass[border = 3pt]{standalone}

\usepackage[svgnames]{xcolor}
\usepackage[3d]{luadraw}

%%%
% La bibliothèque ''shadings'' propre à \tikz nous permet d'avoir
% facilement des dégradés colorés.
%%%
\usetikzlibrary{shadings}

\begin{document}

\begin{luadraw*}{name = gamma-modulus-surface-smooth}
local palext = require 'luadraw_palettes'

local pal = palext.ShiftRainbow

local ld = luadraw

local M   = ld.pt3d.M
local cpx = ld.cpx
local Z   = cpx.Z

local vecK = ld.pt3d.vecK

local graphview = ld.graph3d:new{
  window3d = {-4, 5, -5, 5, 0, 6},
  viewdir  = {-120, 60},
  adjust2d = true,
  bbox     = false
}

local gamma = function(x, y)
  local k_facto, sign, z = 1, 1, Z(x, y)

  local S = 1 / z

  for k = 1, 50 do
    sign    = -sign
    k_facto = k_facto*k

    S = S + sign / k_facto / (z + k)
  end

  return S + ld.int(
    function(t)
      return cpx.exp((z - 1)*math.log(t) - t)
    end,
    1, 10
  )
end

local colors        = {}
local memo_cpx_args = {}
local key

local xykey = function(x, y)
  return x .. "/" .. y
end

local gamma_modulus_pt3d = function(x, y)
  local z = gamma(x, y)
  local arg = cpx.arg(z)

  key = xykey(x, y)

  memo_cpx_args[key] = arg -- save argument

  colors[key] = ld.palette(pal, (arg + math.pi)/(2*math.pi))

  return M(x, y, cpx.abs(z))
end

local color_coef = function(facet)
  local miniarg = math.huge

  for _ , A in ipairs(facet) do
    key = xykey(A.x, A.y)

    local arg = memo_cpx_args[key]

    if arg ~= nil then
      if miniarg > arg then
        miniarg = arg
      end
    end
  end

  return (miniarg + math.pi) / (2*math.pi)
end

local cpx_surf = ld.surface(
  gamma_modulus_pt3d,
  -4, 5, -5, 5,
  {50, 50}
)

cpx_surf = ld.cutfacet(
  cpx_surf,
  {6*vecK, -vecK}
)

cpx_surf = graphview:Sortfacet(cpx_surf)

graphview:Dboxaxes3d({
  grid      = true,
  gridcolor = "gray",
  fillcolor = "LightGray"
})

------
-- Les nouveautés commencent ici.
------
graphview:Lineoptions(nil, "gray", 1)

for _, facet in ipairs(cpx_surf) do
  local facet_sorted = table.copy(facet)

------
-- Nous ordonnons les sommets de façon lexicographique avec
-- l'abscisse prédominant sur l'ordonnée. Ceci va permettre
-- de définir correctement le dégradé de la face.
------
  table.sort(
    facet_sorted,
    function(A, B)
      return (
        (A.x < B.x)
        or
        ((A.x == B.x) and (A.y < B.y))
      )
    end
  )

------
-- Construction des couleurs extrémales qui vont définir le
-- dégradé de la face.
------
  grad_cols = {}

  for _ , A in ipairs(facet_sorted) do
    key = xykey(A.x, A.y)

    local col = colors[key]

    if col ==  nil then
      gamma_modulus_pt3d(A.x, A.y)

      col = colors[key]

      if col ==  nil then
        col = ld.palette(pal, color_coef(facet))
      end
    end

    table.insert(grad_cols, col)
  end

------
-- ¨Def du dégradé ¨tikz morceau par morceau en ¨fonc du ¨nbre
-- de couleurs obtenues précédemment.
------
  local options = "upper left = " .. grad_cols[1]

  if #grad_cols > 1 then
    options = options .. ", upper right = " .. grad_cols[2]
  end

  if #grad_cols > 2 then
    options = options .. ", lower left = " .. grad_cols[3]
  end

  if #grad_cols > 3 then
    options = options .. ", lower right = " .. grad_cols[4]
  end

------
-- Nous pouvons ajouter notre facette "dégradée".
------
  graphview:Dpolyline3d(facet, true, options)
end

graphview:Show()
\end{luadraw*}

\end{document}
